% ============================================================
%  Poul Martin Møller,
%  "Preparations toward a Treatise on Affectation"
%  English translation of "Forberedelser til en Afhandling om
%  Affectation" (Efterladte Skrifter, bd. 3; printed pp. 163--188).
%
%  Mirrors transcription.tex 1:1. \opage{N} marks the point at
%  which printed page N begins. Danish quotes „…" -> ``…''.
%  The connected opening treatment (1837) is followed by the
%  editor's thematic grouping of the author's scattered notes;
%  group-divisions are marked by a centred rule (\ombreak).
%  "Sædelighed" is rendered "ethical life"; "Moralitet" "morality".
% ============================================================

\documentclass[12pt]{article}
\usepackage[utf8]{inputenc}
\usepackage[T1]{fontenc}
\usepackage{libertinus}
\usepackage{libertinust1math}
\usepackage{textalpha}
\usepackage[english]{babel}
\usepackage[protrusion=true,expansion=true]{microtype}
\usepackage[margin=1.4in]{geometry}
\usepackage{setspace}
\usepackage{fancyhdr}
\usepackage{marginnote}
\usepackage[colorlinks=true,linkcolor=black,urlcolor=black,
  pdftitle={Preparations toward a Treatise on Affectation},
  pdfauthor={Poul Martin Møller}]{hyperref}
\setlength{\headheight}{15pt}
\pagestyle{fancy}
\fancyhf{}
\fancyhead[LE,RO]{\thepage}
\fancyhead[RE]{\textit{On Affectation}}
\fancyhead[LO]{\textit{Poul Martin Møller}}
\setlength{\parindent}{1.5em}
\setlength{\parskip}{0pt}
\onehalfspacing

\newcommand{\opage}[1]{\marginnote{\footnotesize\textit{[#1]}}}
\newcommand{\ombreak}{\par\medskip\begin{center}\rule{0.32\textwidth}{0.4pt}\end{center}\medskip}

\title{\textbf{Preparations toward a Treatise\\[0.2em] on Affectation}\footnote{For a long time Poul Møller had had his attention directed to this side of life and had written down his remarks upon it, to which there attached itself the intention of working up this subject in a separate work. In reading the preceding »Strøtanker« [Stray Thoughts] one will have noticed this standing direction in his observation, so that on almost every page traits of falsity and self-deception in the life of men are to be found. Among these were also the notes that here follow separately, and which are chiefly such as the author himself had marked as written down with reference to that concept. This separation is arbitrary, since much of what has been left standing in the collection belongs here just as well; but if everything that could be referred to the concept of affectation had been removed from the main collection, the latter would have lost its peculiar, in a way biographical, character, and some arbitrariness would still remain in determining the limit up to which, according to one's own interpretation of the tendency of the written thoughts, it was permissible to go. Moreover the author has only loosely indicated the manner in which he would have connected this numerous mass of remarks and particular traits, and far from collecting them to fill out a preconceived scheme, he strove especially to draw out affectation in its more latent expressions; so that these observations, as they are had from the author's hand, would in any case work only rather sporadically upon the reader. Those which now follow have been arranged by the editor in groups according to their content. In the year 1837 Poul Møller began a connected scientific treatment of this subject, and what he thus wrote down is communicated here first.}}
\author{Poul Martin Møller}
\date{Posthumous Writings, vol.\ 3 (Fragment IV)}

\begin{document}
\maketitle

\begin{center}\rule{0.32\textwidth}{0.4pt}\end{center}
\medskip

\opage{163}In the reflections here communicated the author has indeed sought to make clear the concept of affectation, but that is --- as a not \opage{164}wholly cursory reader will himself easily become aware --- not what has constituted his main interest. It has mattered to him most to go through the most common phenomena of affectation, both in human life in general and in particular in the age. He who finds no taste in this kind of moral natural-description is therefore warned in good time to skip these pages.

Those who set no store by any thinking other than that in which all concepts are, by an immanent development, produced out of nothing, will naturally demand that the concept here treated shall come forth as the result of a speculative system set out in advance. But if these have not, in consequence of the warning already given, ceased to follow us, they are requested to drop their demand, since it is here held to be unfitting to give the ground-plan of a temple in order to find in it a hole for a poor church-mouse. If all their theoretical and practical concepts really do make up one system, then they will surely have determined therein the concept here treated with the same limits as shall be indicated in what follows. --- The other class of readers, whose several concepts each lead a more independent existence and have more arranged themselves in the form of an archipelago, will probably find room for these reflections in one group or another. In the opposite case, these may well produce that little agitation in the ocean of their thoughts by which a small islet can come forth and subsist for half an hour's time.

Various profound persons in our time make use, in writings and conversation, of the word truth, which they seem to take in such a sense that they thereby denote almost everything that is estimable in human life. On the other side they then condemn as untruth, or indeed even as inner lie, so many kinds of states of mind conflicting with reason, that one is naturally brought to raise the question whether truth, with \opage{165}them, is merely one of the essential sides of ethical life, or whether all the determinations of ethical life are united there, or whether it is a peculiar virtue alongside other virtues? This the many who make use of that concept in aphoristic utterances seem not really to have made clear to themselves; and there is hardly any systematic work in which that view is so carried through that the indeterminate concept which hovers before those persons can be assumed to have its adequate presentation therein. From the short registers one sometimes finds of the various moral principles that have been asserted in the history of practical philosophy, one might be tempted to believe that the Englishman Wollaston had precisely conceived moral perfection as truth; but on reading his own work one does not find in it the view here spoken of. Nor is any working-out of it found in Ammon's moral system, although in looking through his chapter on the moral principle one will be brought to believe it.

A man's life has a kind of truth when, without dissimulation, it follows its natural desire: a higher truth it has when it has attained virtue (in the ancient sense of the word), so that it still draws the content of its actions from the natural drives, but has yet won such a mastery over them that in their satisfaction it observes a certain measure. A still higher stage of personal truth is had by the life of him who, with pure rational autonomy, determines all his resolves. (That this cannot happen by a merely subjective thinking, but in such a way that the subject recognises the rational order, which without its co-operation is in its development in existence, to be a work of the same reason that is the proper truth of its will --- this does not belong here to be proved.) In so far, finally, as man's pure self-determination is the will hallowed by religiousness, he acts in complete harmony with the whole world of reason; he is what he shall be, and his life can \opage{166}attain no higher truth. But this truth is nothing but ethical life, and every deviation from it is unethical conduct. Affectation is certainly a species of this untruth; but we have to determine more closely what place it occupies in this larger circle.

The man who has affectation in his life determines himself, in so far, not with complete moral freedom; his actions do not have their origin in the true self, which is his free ethical will. His will is determined by some merely natural aim or other, whereby he is brought to represent an alien person or to take upon himself a false role that is not assigned him in life.

It belongs then, in the first place, to the essence of affectation that it is a falsity. But not every falsity is affectation. He who, with clear consciousness that he is lying or dissembling, makes himself guilty of lie and dissimulation, does not thereby lay any affectation to the day. We by no means say hereby that affectation is any lesser vice than the dissimulation that is accompanied by clear consciousness. It will appear in what follows that one degree of affectation presupposes less immorality than deliberate dissimulation, while in another degree of it a greater immorality is laid to the day. Affectation, then, is not unmixed falsity, but always has an admixture of self-deception; for it lies in the concept of affectation that the person strives to be what he cannot be; but he cannot strive after this without at least for a while imagining to himself that he can be it. But the wholly inculpable self-deception, which is without any falsity, can just as little be called affectation; so that the two constituents here indicated, the falsity and the self-deception, are always found combined in it. (In order to attain an aim lying outside the action.)

(It ought to be a superfluous remark that wholly invol\opage{167}untary movements can impossibly be reckoned as affectation, and that consequently the abnormal expressions of life which, owing to the spiritual life's lack of command over the bodily, not seldom escape a man wholly against his will, must be excluded from it. It is nonetheless not seldom that one hears such unnatural and inadequate expressions of the inner life labelled with the name of affectation. We think here, in the first place, of all the outward movements which, on account of organic defects, do not stand in that relation to the person's representations, thoughts and feelings in which they would stand if the body were in its normal condition. Examples of this are awkward movement of single limbs, involuntary muscular twitchings, faults of pronunciation etc., which together are reckoned by the mob as affectation. Under the same heading belong a number of failed expressions of life which are of just as innocent a kind, although even sensible people not seldom pass a condemnatory judgement upon them. To these one may reckon every disproportionate expression of feelings which becomes disproportionate because the person, either constantly or under certain circumstances, lacks the capacity to find the fitting expression for his mood. Such mistakes can just as little be reckoned as moral faults as the foreigner's words when, from lack of fluency in the language, he comes to make use of unbecoming or offensive expressions.\footnote{In an older draft the following stands here: If, namely, one regards every outward means of communication as a language --- and in this connection mien and movements are indeed just as much a language as the audibly uttered word --- then there are a great many involuntary faults of speech which are without ground reckoned as affectation.} A man of low station, not used to being in company with distinguished people, often makes his bows too deep; not because he wishes to lay so great a reverence to the day, but because he does not have the usual skill in calculating the angle of inclination \opage{168}with accuracy. --- He who, during a conversation, notices that he is coming, against his will, to express himself too harshly and strongly against another, will sometimes make his fault good again, but utters --- because his blood has now once been thrown out of equilibrium --- a benevolence that is not quite so meant, although the miscarried correction may nonetheless have its origin in a striving to lay his true heart's disposition to the day.\footnote{In an older draft: It hardly needs to be remarked that in the depiction of these concepts we start from the inexorable presupposition that the party concerned has no other aim than to express what really dwells within him. If there are other by-purposes involved in the game, as e.g.\ with him who, because he wishes to conceal the slow circulation of his thoughts and his lack of presence of mind, seizes a false expression of his personality precisely because he grasps at it at random, then naturally no full sincerity occurs, although others just cannot judge to what degree it is lacking.} Everything in a man's expressions of life which thus does not at all lie in the subject's preceding representations of them cannot be called affectation. Man in such cases only succumbs to a conflict with the necessity of nature, in that outward organs refuse to obey the commands of his free will. (Ground in preceding affectation\footnote{In an older draft: But hereby it shall by no means be said that the disproportion here spoken of between the personality and its expression cannot have its ground in preceding affectation, (an originally assumed manner, like that of one who takes pleasure in others' wholly individual bodily movements, their gait or voice, either because these expressions considered for themselves appear to him to become the other well, or because the other's distinguished personality so enchants and blinds him that for him it casts a lustre upon such accidental trifles. Thus Plutarch relates that the courtiers of Alexander the Great acquired the habit of letting their head incline a little to one side, because they constantly saw their prince in this posture. Such expressions of life, which have first sprung from affectation, may at last have become a second nature, so that they are no longer called forth by design or accompanied by consciousness.)}).) The assumed role may be imitation which \opage{169}conflicts with one's nature, or a character invented according to one's own liking, which seems to become one.)

If we consider the degrees of affectation according to their greater or lesser connection with the subject's character, then it has the least possible share therein when it conflicts with the direction of the will, but in the single moment escapes the man because his virtue has not become a facility. Such an expression of affectation stands in the same relation to affectation that has become a habit as an intoxication does to habitual drunkenness, and there is hardly any mortal who does not, at least in youth, sometimes make himself guilty of small, half-involuntary approaches of this kind. If there is anyone who, even at the beginning of his period of development, so preserves his independence that he even avoids that shape of affectation which lies near the border of amiability --- that he, e.g., never, out of a desire to be in complete agreement with a friend, forces himself to sympathise more with them than he can --- then he must certainly by nature have a very strongly marked individuality; but whether he therefore has the happiest natural disposition is another question. His gift for keeping his individuality pure and free of extraneous admixtures may also have its ground in an egoistic bent to enclose himself merely in his own circle of thought, and a lack of capacity to open his mind to alien influence. He who is not able, in intercourse, so to give himself over to others that he for a while becomes one with them --- that he goes quite out of himself and loses himself in an alien circle of consciousness --- can certainly by his reserve save himself from being overwhelmed by any spiritual power; but the \opage{170}individuality that is salvaged only in this way always becomes one-sided and poor.

Spiritual perfection can, like physical growth, be furthered only by the individual's regularly melting together with what is alien to it, and apparently sacrificing itself in order to return home enriched to itself again. But then it may certainly happen that a well-endowed youth, from such a fruitful self-forgetting, in which he has let another dispose over his train of thought, returns to himself with the semblance of representations and feelings which are repudiated by his true self, but which he forces upon himself in order to enjoy the ravishing pleasure of sympathy. The expressions of affectation that come forth in that way may often precisely become the ground of a strengthening of the personal life's truth, in that he who has heedlessly overstepped the border of his personality is, out of dissatisfaction with himself, brought to a clearer consciousness of what he feels, wills and knows, and to a more vigorous self-movement in similar cases. He who has been on the far side of the border now first knows where the border is. --- Such a passing expression of affectation, which has its source in a power that blinds feeling, from which the man by a subsequent reflection tears himself loose, as from something that does not belong to his essence, and against which he guards himself for the future, is a lesser fault than the lie proper. (\emph{The momentary affectation.})

A higher degree of affectation occurs in him who has assumed the habit of a definite kind of false expressions, in that he imagines himself to have certain opinions, interests or inclinations because from some outward ground or other he wishes to have them. As when one, out of vanity, lies himself into a love for some art or other for which he has no sense; or when relatives and friends of a zealot count, both to themselves and to others, as being like-minded with him, although his zeal is alien to their heart and under altered circum\opage{171}stances falls away of itself. The possibility of the self-deception here rests on the logical consistency with which the assumed role is pursued; but the centre that bears such a circle of representations falls outside the subject himself. This second degree of affectation is in a moral respect not at all less imputable than the lie and dissimulation accompanied by clear consciousness; for the fact that it is not brought to full consciousness comes merely from the party concerned not willing to bring it to consciousness. This is remarked because a most incorrect concept of it has come into circulation, the multitude believing that in self-deception there lies sufficient justification, or at least excuse, even for open badness --- an opinion whose incorrectness is easily evident when one pursues it into some of its undeniable consequences. From it would namely follow that one could drive one's falsity so far that it again became honesty, and that he who lies so long that he himself believes his invented tales would at this stage become better, because his lie is transformed into self-deception. According to this view a judge could also, in all innocence, let himself be bribed, since he would merely need, by sophisms, to bring himself to believe that his benefactor was in the right. The conviction that has its origin in inclination cannot be particularly well-grounded, which already shows itself from the fact that the effect in general does not last longer than its cause. It is not improbable that he who finds a literary work bad because he believes his enemy has written it will find it very good if he learns that his friend is the author of it. But that in the first case he found it bad was certainly no innocent error, however much self-deception there might be in it; for one ought just as little in one's assumptions as in one's actions to follow one's inclinations against reason and truth. Under this second degree of affectation the man takes a false element into himself and distorts his personality, so \opage{172}that its expressions do not hang together with its proper self.

In so far as it thus has a double inner --- one real, which is thrust back, and one apparent, which it wills shall count for itself and others --- it leads, with respect to the latter, only a life of mere show. (\emph{The fixed affectation.})

The affectation last spoken of, even where it is driven furthest, still allows the person to have a lasting semblance of inner connection, both for the subject himself and for others. But such a semblance of consistency can be only momentary in the third and worst degree of affectation. It occurs where the man does not have some feigned trait or other in his character, where he does not have a habit toward a definite kind of affectation, but has a facility in affectation in general, which now assumes this, now that definite shape. This badness approaches, according to its greater or lesser development, more or less to complete untruth in the personal life. If a man could reach its culminating point, then there would be no abiding kernel in his thinking and willing, but there would form itself, in every moment of his life, a temporary personality, in order to be abolished in the next. He would probably for the most part, like certain animals, change colour according to his surroundings, and in so far be the passive product of his circumstances; but since this is only one of the shapes of affectation, his course could nevertheless not let itself be calculated according to the simple rule that he should resemble his surroundings, since affectation can also show itself in a striving to present, in one's conduct, the peculiar, the unusual. To this completed lie in the inner life no one can bring it; but if anyone could, then it would be a moral suicide, whereby he had wholly annihilated himself as a peculiar figure in the moral world. (\emph{The changing affectation.})

After we have in the foregoing developed the three main species of affectation with respect to their greater or lesser connec\opage{173}tion with the character, we shall now consider the different ways in which it expresses itself, and the different grounds from which it takes its origin.

\ombreak

Affectation can show itself in this, that one with half-design imitates the peculiar in others' involuntary expressions of life, because it seems to become one well. It is a kind of affectation which is sometimes a pure fruit of love and admiration, without any other aim. The assumed, alien element in the movements which are supposed to be involuntary, and which no one can strive to alter without there being some curable natural defect in it, is called by all most of all affectation, although it is the most insignificant kind of affectation. --- To imitate the accidental in others' speech. To use philosophical terminology without the philosophical insight.

It can lay itself to the day in feigned feeling. Men then put themselves so vividly in the place of those who really have such a feeling that they imagine to themselves that they have it. That one can do: warm actors' example shows it, since they can so identify themselves with a role that they forget who they are. --- Occasional poets do it \emph{sensu pejore}, who sometimes put themselves into a mournful mood, sometimes into a joyful one, according as one has ordered a funeral poem or a wedding song from them. That people can imagine to themselves that they have feeling comes, then, from their using the language and expression of feeling, and thereby producing in themselves a kind of reminiscence of the true feeling. --- The man who talked himself into grief over the \opage{174}deceased relative.\footnote{Cf.\ p.\ 34.}

One can take upon oneself a self-feeling one does not have, just as well as a social one. Also the higher contemplative feeling lets itself be counterfeited by an art that deceives both others and the affected themselves. Especially is enthusiasm for art a common affectation. (Feigned enthusiasm --- merriment.)

Affectation can show itself in cognition, when one, without conviction, out of inclination assumes propositions to be true.

Affectation also shows itself in this, that the man under self-deception expresses himself as though he had an inclination, a natural calling, which he does not have. --- (Artist.)

\ombreak

Affectation always has its ground in the man's being bribed by some inclination or other, without himself knowing it. For the affected man does, by design, what has no significance, except where it is called forth by natural power and natural prompting or higher necessity. This inclination is now self-interest, now weakness.

Thus self-interest brings many a priest of superstition to assume articles of faith, because his station in life rests on their being true. A Catholic clergyman, e.g., must proclaim many a doctrine which is in itself so conflicting with reason that one cannot suppose that so many could be convinced of its correctness if they did not, for the sake of convenience and advantage, deceive themselves.

Weakness too can be the source of affectation. Men can, out of benevolence, so let themselves be carried away that they force upon themselves a conviction in order to be in agreement with the one they love. How else could it come about that people, in scientific feuds, where conviction should be determined merely by grounds, are yet in general on the same side as their friends, kinsmen and relatives? \opage{175}But this desire to be in agreement with others, and therefore to become so, does not always come from benevolence: it may also have its ground in vanity. This is the other weakness that often occasions affectation. In order to be agreeable to others one often becomes of one mind with them. This kind of affectation is found especially in sanguine persons and can, when they have outstanding mental gifts, produce remarkable phenomena. There are persons with especially quick heads who are to such a degree a product of their surroundings that they wholly take on its colour. A high degree of suppleness in the mind's faculties and receptivity for the most diverse views make them attach themselves now to one, now to another literary or political party, and yet be with life and soul in each of them, so long as it lasts. Such incoherent persons live, namely, always in an immediate cognition, without being able to raise themselves to the standpoint from which different one-sided standpoints can be surveyed and seen in their relative validity. When this kind of turncoat possesses a high degree of mental riches, as they often do, they can in a brilliant manner carry out the person that is given them from without, without having any deep root in their own personality.

It is not merely a vain desire to be included and to find good comradeship that is to blame for such an approach; it can also be a striving after the pleasure that a heightened self-feeling grants. (\emph{Moral reveries}.)\footnote{\textbf{Alibi:} An example of the affectation springing from the desire to please oneself is fantasies over plans for activity which one half-and-half knows one will not carry out.}

Every personality is something infinitely glorious when it is carried out according to its idea; when, on the contrary, one lets it be overwhelmed by alien considerations, one fails one's vocation. Perhaps \opage{176}more dangerous to morality than conscious lie is the untruth mixed with self-deception, whereby the man amalgamates his false person with the real one. The fully conscious lie one can lay aside after a serious resolve; but one can grow fast in the false person without being able to tear oneself loose from it. One makes oneself into a moral changeling.

\ombreak

When one gives oneself over confidently to an idea and follows its guidance without anxiously reconciling every proposition with current views, one is readily held to be fantastical. That I shall certainly be held to be, when in my work on affectation I maintain \emph{that no expression of life has truth unless there lies in it creative self-activity} (p.\ 91).

In the Introduction (on the tyranny of concepts) it must indeed be developed that one does not think in words merely.

There are two kinds of truth in human life. The one is natural, innate: thus an inclination to affectation can be innate. The other is chosen and won with freedom; it is of higher and richer content.

Some self-deception is necessary; for it makes up life.

\ombreak

Affectation comes most often from one's not having the strength to fall out with the world in order to show one's personality true. Therefore it is good that some come to stand in defiant opposition to the guild. In the state of nature each lived separate and developed, in defiance, a personality, since he was not disturbed by the many. Now one forms for oneself, by abstraction, a universally valid \opage{177}person, a social ideal without edges, an ideal without individuality. (Nature then asserts its right in books like Baggesen's Ghost.)

Opposition too can be driven with affectation; one can strengthen the disharmony because one has sacrificed one's life to war.

Everyone has by nature his deep stamp, but lets it be effaced in order to please others. These others have, in rejecting those who protrude, a very correct principle; only they should make a distinction between the self-willed and those who protrude out of pedantry and helplessness.

The wholly dissolved persons, whose life is only mimicry, sometimes have so rich a nature that they can carry out an assumed character with such force that the individuals who really have this character cannot carry it through more forcefully.

The incoherent are too bad to be evil.

\ombreak

Affectation can come from one's imagining that one is working for an idea, and thereby procuring for oneself a heightened self-feeling.

Some ascribe to themselves false qualities merely in order to say something about themselves. --- Hysterical women ---.

Affectation of natural advantages: When a schoolmaster wishes to be a youth's god --- an Apollo. Or when one curses and uses a posteriori witticisms in order to seem a strong spirit, when yet it is really merely indolence that brings one to such a thing. Or \opage{178}when an old maid wishes to be young. Or when a Doctor Faust wishes to be a romantic lover, or an old fool to be jolly with young students; Rahbek to be a bold warrior. --- Between renting and owning life hovers. Youth rages, said the old woman, as she leapt over a straw. Cicero wishes to be a hero of freedom. By his zeal for freedom nothing was achieved; it was merely for show.

One should precisely lay bare great men's faults in order rightly to learn to know human nature. Almost all cultivated people lie in the telling of their love.

It is often a mistake that Many believe themselves able to write great works; he holds for himself a false feast of annunciation.

To coquet with open-heartedness, modesty, etc.; with unconscious abandon, naïveté, a childlike mind. The affectedly frank and open-hearted man narrows his eye a little; the self-satisfaction that shines from his eyes betrays him, and the half-constrained smile that plays about his mouth.\footnote{Cf.\ p.\ 54.}

``The ass could not see that it was precisely childlike of me that I gladly wished to eat pancakes.''

One of the causes that human motives lie so wholly in darkness is the circumstance that people, in order to seem noble, say, on hearing observations about egoism: No, God forbid, that is not true. --- Also, by disapproving of the least spot on tragic heroes, they lay their noble soul to the day. Misanthropes, on the contrary, exaggerate everything. --- (Spleenful Englishmen or French despairing deniers of God.) \opage{179}Those who adorn themselves with morality (which is just as wrong as to adorn oneself with immorality) use ``God pardon him,'' in the same sense and with the same embittered look as when they would say: ``God damn him.''

Servant-girls say, in order to provoke one another: God pardon you your sin, Grete! but you are an evil Satan --- in a meek tone, so that they come themselves to believe themselves noble.

(To surround oneself with a halo.)

That no one easily falls, for a sister or brother, into moral declamations shows that the rigoristic declaimer uses his zeal as a Sunday dress.

How great is not the difference in the conduct of fourteen-year-old girls and boys in the presence of strangers? These quarrel rudely with their mother in strangers' hearing; those behave gracefully --- from an unconscious desire to please.

The girl fondles the cat in order to seem tender-hearted.

It is often, in order to give oneself a semblance of omniscience, that some say: Everything is old.

An affectation in learned men and geniuses is to seem not to be of this world. They hear with joy their wives talk of their distraction, unconcern and ignorance in the affairs of daily life, since they are not of this world.

Consider how much lie can be found in an everyday scene of life. The traveller asks Bernhard: ``How \opage{180}does one make pancakes?'' Lovise: ``That my learned brother does not know.'' Bernhard: ``Oh, yes, indeed I do know'' (gives out feignedly some incorrect ingredients and nudges his neighbour Augusta in the side). She says: ``Yes, shameful to speak of'' (does she mean it?) ``I don't know it either'' (does she speak truly therein?).

Sometimes a whole class plays its role, like S. and distinguished sailors.

Sailors of rank mumble in imitation of the real so-called salts, with half-open mouth, or use a manner of speech as though they had a quid of tobacco in the mouth, though having none.

\ombreak

It is a droll kind of weakness in certain people that politeness tyrannises them. --- \textbf{A.} ``You absolutely must drink coffee at Mini's, it is better than Semadeni's.'' \textbf{B.} ``No; I would rather go to Semadeni's.'' --- \textbf{A.} ``I assure you, you absolutely must drink your coffee at Mini's.'' --- \textbf{B.} ``No'' --- \textbf{A.} (pushing him along) ``Yes, you absolutely must.'' --- \textbf{B.} pretends he will go in, and lurks in the gateway until \textbf{A.} is some way off; then he goes his way. The one's importunate politeness and the other's nightcap-like delicacy are equally remarkable. So touched and constrained do people become out of politeness and solicitude for those whom, on both sides, they know to be a semblance.

It is a very innocent kind of affectation to set the muscles in the position they take when one laughs or smiles, while one hears an anecdote which does not amuse but nonetheless lays claim to it. \opage{181}He who is glad at compliments during disputations knows nonetheless very well that those who make them do not mean them. He deceives himself voluntarily in order to enjoy a little triumph. --- The Great too continue to believe in flattery.

Praise and admiration become in time a physical necessity for old men. They wish to live in a dream of having been Something, although they are Nothing; they wish to deceive themselves.

Reconciliation is often self-deception. ``Now I have made peace with him, the ass!''

Is it not affectation when a proud choleric wishes to conceal that he wonders at something or finds it ridiculous etc., but expresses that he does not find it so?

Affectation arises sometimes from haughtiness. He who wishes to conceal the effect that others' talk makes upon him is at once untrue and caught in self-deception. The latter, because he is not, in the same moment as the untruth is expressed, conscious of it, although by a subsequent reflection he often must be able to recognise his falsity.

\emph{To put on airs.}

Dread of affectation can --- however strange it may at first glance seem --- be driven so far that an abnormal state of mind thereby comes forth. This dread produces a misunderstood holding to abstract identity with oneself. Thus I have known a young man who could not persuade himself to send off a letter that had lain a day over in his desk, because it seemed to him to be untrue, since it namely no longer expressed his present mood.

\ombreak

\opage{182}An affectation is the lifeless echoing of the philosophical folk-legends about the glory of art. Not merely art in the narrower sense, but the art of life deserves prize. --- Is composure not enthusiasm? then there is no enthusiasm in nature's calm rational creating, and it is yet just that which is to be imitated. --- Dizziness --- the blood to the head --- drunken derangement, and the so-called inspiration, which is really a deranged gladness over the blood's seething.\footnote{The momentary exaltation is often of a low sensuous nature. Author's Note. Cf.\ p.\ 89, third aphorism.} --- It is often vanity's rapture over single felicitous expressions. Enthusiasm is the calm working of a model laid down in us by nature. It can be better imitated by calm diligence than by a rushing in the ears. There should be pleasure only in the feeling of the whole's felicity, not in the single word and expression. Does the calm chiselling-out of a Greek statue not presuppose enthusiasm? It is an ever-recurring affectation to conceal one's diligence.

In Homer's time all speech was treated as a work of art. --- Homer is not written in a drunken frenzy.

A lie is the poetry that does not come from life. The nearer it comes to life, the more true; the more distant, the more lie.

The moral fantast lacks fantasy; for his phantasms are abstract.

The present generation has so high a degree of excitability that it is set in emotion by demonstrating the objectionableness of emotion with respect to artistic or scientific production. By the vivid representation of the loftiness that lies in the preponderance of cognition and will, and the unhealthy \opage{183}direction that is laid to the day in the current sentimentality, it finds nourishment for its own sentimentality. The most dissolved in feeling, who would necessarily have to come to consciousness of this their sentimentality, strive with the greatest pathos against sentimentality.

\ombreak

The will shall not determine what is true.

Spurious enthusiasm can express itself in this, that one by vehement declamation wishes to make one's will prevail, without procuring conviction by grounds.

Many will not part with their errors, which with long entangled roots have, like stubborn weeds, laced themselves fast about a great many neighbouring plants. They believe they are fighting for truth and fight for their own persons.

It is always affectation when a judgement does not express a thought that is determined by an object, but partly by an inclination. He who has it in mind to write regards himself as belonging to the party of writers and is often forbearing, but becomes exaggeratedly severe when he has resigned.

People shut their eyes to the Greeks' shadow-side and to tyrants' light-side.\footnote{See p.\ 37.} Affectation it is when one who has for a long time shown an inclination to bring out the shadow-side in everything suddenly wishes to show himself willing to recognise the better side of things and suddenly falls to praising this and that. --- Likewise when people, out of fear of showing sensitivity, pretend to be hard-hearted, --- and \opage{184}when one forces oneself to acknowledge distinguished people's merits, in order not, in one's own and others' eyes, to pass for a flatterer.

A kind of affectation is the unphilosophical multitude's contemptuous utterances about a philosophy that is above their horizon. They set the limits of their own cognition, but with weak conviction, as the limits of reason. (He who cannot understand philosophy prefers that there be nothing in it.)

\textbf{Par force} wishing to favour some trash out of fear of one's own criticism. --- It is an affectation to deny one's critical nature.

When the boy says: I think more of Goethe than of Schiller, or in general when the immature think more of objective than of subjective poetry, then it is as when the child denies itself the loving of sweets and forces itself, before its time, to love the bitter.

Affectation is Tegnér's talk of the fusion of form and matter among the Greeks, which he himself contradicts by delivering preponderant form.

\ombreak

Does he speak truth who speaks words he does not himself understand? --- No.\footnote{Cf.\ p.\ 30.} Does he then speak truth who makes use of words whose concept has been imparted to him by another, without their standing in any organic connection with his remaining cognition? Does he speak truth who disputes with words he does not understand, who is governed by his own definitions? --- But it is indeed merely \opage{185}word-play; he who believes he speaks truth is \textit{bona fide}. --- No, he deceives himself. He is ashamed to stand naked, in that he casts away all traditional concepts, which make up his whole scientific store. He will not come forth with his own person, does not believe in his person's infinite depth. If every single man, without fearing censure for simplicity, judged of things as they presented themselves to him, then glorious phenomena would come forth. But most fear to lay mental poverty to the day. They confirm confused heads in their confusion by forcing themselves into these people's views, out of fear of being taken for people of poor understanding. --- Therefore Danish people force themselves to speak Germanically; therefore people esteem philosophical phrases which they do not understand.

He does not speak truth who says that man does not give himself laws; for he does not know what an I is.

An affectation is morality according to dead concepts. He who, e.g., says: ``Fie! A is a suicide,'' has only a word to hold to. He who says: ``A ought not to have forsaken his old father,'' speaks from intuition, immediate cognition.

A peculiar kind of affectation is that of an eclectic when he uses systematic philosophers' terms; which is unnecessary for him, whose inorganic philosophy can be understood at every point.

People hold to distinction in the general concepts.

\ombreak

A peculiar kind of affectation in discourse is the grammatical. This causes lexical studies and care for \opage{186}a schooled syntax to step visibly forth in the style. When one reads such writings, it is as though one read a dead language. Such purists set up a factory for native words which are to replace other current ones whose citizenship is suspect. These home-made words they keep in their savings-box, and when they have collected a fitting number, they take occasion from it to write on some trivial subject or other, which allows them to grant undivided attention to the word-construction. When one rightly considers their products, the language is the real content, and the apparent content a by-thing.

Affectation: to write in a decisive tone, violently to close one's account-book of cognition. To that one forces children by having them write Danish compositions. One makes them into stylists, which is one of the worst things a man can be.

It is strange that people do not lose the desire to be composition-writers by seeing with what loathing they read others' compositions. Properly compositions should be written only with negative satisfaction, with the conviction that one used fitting expressions.

Pleasure in using certain words: Dannemand [man of worth], well-nigh.

A proof that gladness over style is a lie is that a boy blushes when one makes him aware that he uses ``tarry'' for ``linger,'' ``rejoinder'' for ``answer,'' or declaims about silver-white old men.\footnote{On stylism cf.\ pp.\ 131, 35, 23.}

\ombreak

\opage{187}There was once a man who seldom went out to walk, although his bodily constitution made it desirable that he had done so. This lack of movement was, on the contrary, replaced in another way by a natural instinct; he talked himself, namely, into an affect over some matter or other, and his blood was thereby set in such a circulation that he had need of no other motion. Since every matter can be seen from two sides, it was accidental from which side it was momentarily seen. The momentary clear view of an isolated proposition's correctness, taken in the subjective sense in which he just accidentally took it, served him as a vehicle for a tumult in the circle of representations, which soon again subsided, without leaving remarkable consequences in it. The author of this book once saw in China a bundle of frogs which were held for sale in a market-place, and which in general lay without remarkable movement, because the animals tied together strained their forces in such opposite directions that the result of their exertions became about nil. Sometimes, on the contrary, a single frog succeeded, without resistance, by a bold leap in lifting the whole bundle for a moment into the air. So it went with the said man's representations; sometimes a single representation became so dominant that his discourse took on a passing semblance of systematic connection.

The scientific workman collects mechanically what is wont to stand together in a discipline, without cognition of its belonging-together.

An affectation it is that he who wishes to write a treatise thinks of a title for himself and is confined by it.

It is a kind of affectation to make a treatise complete because single chapters are alive for us. \opage{188}That the Ancients did not have the systematic affectation is seen from their frequent digressions. They were not tyrannised by their titles. Everything that is life strives against stiff forms.

Affectation is symmetry in philosophical systems.

\ombreak

It furthers love of humanity to make clear to oneself concepts of bad qualities, in order that they may be sharply distinguished from the good. (Some find egoism in everything.)

A beginning judge of men's insight into badness is self-deception. Guard yourself against coming into his neighbourhood. He begins, like the toddler who wishes to learn to know a thing, by tearing it apart. But all conjecture of badness or gloriousness vanishes in the neighbourhood.

Insufferable judges of men, who would force people to act according to their representations, and craven wretches, who dare not set themselves against the others' concepts!\footnote{Cf.\ pp.\ 21, 26.}

\ombreak

\end{document}
